% Weak Independence Theory

\vspace{3pt}
\noindent Signal hierarchy theory (Section 4) enables unified metabolic-regulatory models by formalizing information channels with consumption semantics. However, such unified models integrate hundreds of metabolic reactions exhibiting convergent coupling (shared outputs like ATP, NADH) and regulatory coupling (shared catalysts like allosteric enzymes). Classical Petri net strong independence prohibits these coupling modes, forcing sequential execution and making large-scale unified models computationally intractable. To provide the computational foundation necessary for signal hierarchy, we formalize weak independence theory distinguishing competitive from non-competitive coupling and proving correctness for parallel execution. This enables unified models where metabolic concurrency and regulatory hierarchy coexist.

\subsection{Motivation and Intuition}

Consider two reactions producing ATP via distinct pathways:
\begin{align*}
\text{Glycolysis:} & \quad \text{Glucose} \to \cdots \to 2\,\text{ATP} \\
\text{Oxidative:} & \quad \text{NADH} \to \cdots \to 3\,\text{ATP}
\end{align*}

Both share ATP as output (convergent coupling). Rates superpose according to ODE superposition principle:
\[
\frac{dM(\text{ATP})}{dt} = r_{\text{glycolysis}} + r_{\text{oxidative}}
\]

\vspace{3pt}
\noindent No substrate competition exists---input places are disjoint. Reactions are genuinely independent at the conflict level. However, classical strong independence prohibits shared outputs, forcing sequential execution despite biological concurrency.

\subsection{Formal Definitions}

\begin{definition}[Strong Independence]
Transitions $t_1, t_2 \in T$ are \textit{strongly independent} if:
\[
({}^\bullet t_1 \cup t_1^\bullet \cup \Sigma(t_1)) \cap ({}^\bullet t_2 \cup t_2^\bullet \cup \Sigma(t_2)) = \emptyset
\]
where ${}^\bullet t$ and $t^\bullet$ denote input and output places respectively, and $\Sigma(t)$ denotes regulatory places connected via test/inhibitor arcs.
\end{definition}

\vspace{3pt}
\begin{definition}[Weak Independence]
Transitions $t_1, t_2 \in T$ are \textit{weakly independent} if:
\[
{}^\bullet t_1 \cap {}^\bullet t_2 = \emptyset
\]
\textit{i.e.}, no shared input places (substrates).
\end{definition}

\vspace{3pt}
\noindent Weak independence permits two non-competitive coupling modes: \textbf{Convergent coupling}---$t_1^\bullet \cap t_2^\bullet \neq \emptyset$ (shared outputs). Rates superpose without conflict. \textbf{Regulatory coupling}---$\Sigma(t_1) \cap \Sigma(t_2) \neq \emptyset$ (shared test/inhibitor). Both read state simultaneously.

\vspace{3pt}
\noindent\textbf{Stochastic parallelism}: Approximate stochastic simulation methods decouple reaction events within time intervals, enabling concurrent Poisson sampling for weakly independent transitions. This reflects biological reality---spatially distributed molecular collisions occur in parallel, not sequentially.

\vspace{3pt}
\begin{definition}[Competitive Coupling]
Transitions $t_1, t_2 \in T$ are \textit{competitive} if:
\[
{}^\bullet t_1 \cap {}^\bullet t_2 \neq \emptyset
\]
Shared input places create true conflicts requiring mutual exclusion.
\end{definition}

\subsection{Correctness Proof}

\begin{theorem}[Safe Parallel Integration]
For continuous transitions $t_1, t_2$ that are weakly independent, parallel ODE integration preserves correctness.
\end{theorem}

\begin{proof}[Proof sketch]
Let $r_1 = \Phi(t_1)$ and $r_2 = \Phi(t_2)$ denote reaction rates. Consider three cases:

\vspace{3pt}
\noindent\textbf{Case 1: Convergent coupling} ($t_1^\bullet \cap t_2^\bullet \neq \emptyset$). For shared output place $p$:
\[
\frac{dM(p)}{dt} = r_1 \cdot \text{Post}(t_1, p) + r_2 \cdot \text{Post}(t_2, p)
\]
By ODE superposition principle:
\[
\int_0^{\Delta t} (r_1 + r_2) \, dt = \int_0^{\Delta t} r_1 \, dt + \int_0^{\Delta t} r_2 \, dt
\]
Parallel integration yields identical result to sequential.

\vspace{3pt}
\noindent\textbf{Case 2: Regulatory coupling} ($\Sigma(t_1) \cap \Sigma(t_2) \neq \emptyset$). Test/inhibitor arcs read marking without modification. Both transitions evaluate regulatory function $\text{Reg}(p)$ simultaneously on identical state $M(p)$. No race condition exists.

\vspace{3pt}
\noindent\textbf{Case 3: Disjoint inputs} (${}^\bullet t_1 \cap {}^\bullet t_2 = \emptyset$). Substrate consumption occurs in disjoint place sets. No interference possible.

\vspace{3pt}
\noindent Therefore, parallel execution preserves correctness for all weakly independent transitions.
\end{proof}

\subsection{Dependency Classification Algorithm}

We provide efficient algorithm for pairwise classification:

\begin{algorithm}[h]
\small
\caption{Classify Dependency}
\begin{algorithmic}[1]
\Procedure{ClassifyPair}{$t_1, t_2$}
    \State $I_1 \gets$ input places of $t_1$
    \State $O_1 \gets$ output places of $t_1$
    \State $R_1 \gets$ regulatory places of $t_1$
    \State \textit{Similarly for} $t_2$
    \If{$(I_1 \cup O_1 \cup R_1) \cap (I_2 \cup O_2 \cup R_2) = \emptyset$}
        \State \Return \textsc{StrongIndep}
    \ElsIf{$I_1 \cap I_2 = \emptyset$}
        \State \Return \textsc{WeakIndep}
    \Else
        \State \Return \textsc{Competitive}
    \EndIf
\EndProcedure
\end{algorithmic}
\end{algorithm}

\textbf{Complexity}: $O(|F|)$ per pair, $O(|T|^2 |F|)$ total for all pairs.

\subsection{Parallel Execution Strategy}

Weakly independent transitions can integrate in parallel via: \textbf{(1) Dependency graph construction}---Classify all transition pairs. Build directed acyclic graph (DAG) where edges represent competitive dependencies. \textbf{(2) Concurrent integration}---Transitions without incoming edges execute in parallel. Use thread pool with $k$ workers where $k$ is number of CPU cores. \textbf{(3) Convergent synchronization}---Accumulate contributions to shared output places using atomic operations or synchronization barriers. \textbf{(4) Regulatory reads}---Test arc evaluations occur without locks (read-only access).

Expected speedup scales with degree of weak independence in the network.
